\documentclass[../../main.tex]{subfiles}
\input{0-theme/default}

\begin{document}

% ==========================
\section{Gasificación}
% ==========================

%%%%%%%%%%%%%%%%%%%%%%%%%%%%%%%%%%%%%%%%%%%%%%%%%%%%%%%%%%%%%%%%%
\subsection{Introducción}
\begin{frame}{Introducción}


\begin{itemize}
  \item El diseño y operación de un gasificador requiere comprender:
  \begin{itemize}
    \item El proceso de gasificación.
    \item Configuración, tamaño y tipo de materia prima.
    \item Parámetros operativos y su impacto en el rendimiento.
  \end{itemize}

  \item Conocer las reacciones básicas es esencial para:
  \begin{itemize}
    \item Planificación, diseño y operación del sistema.
    \item Diagnóstico de fallas y mejora del proceso.
  \end{itemize}
\end{itemize}
\end{frame}

\begin{frame}{Gasificación: Reacciones y Etapas}
\small
\begin{itemize}
  \item La \textbf{gasificación} convierte materiales sólidos o líquidos en:
  \begin{itemize}
    \item Combustibles gaseosos útiles.
    \item Materias primas químicas para productos con valor agregado.
  \end{itemize}
  \item \textbf{Comparación con la combustión:}
  \begin{itemize}
    \item Ambos son procesos termoquímicos.
    \item \textbf{Gasificación:} almacena energía en los enlaces químicos del gas producto.
    \item \textbf{Combustión:} rompe enlaces para liberar energía.
  \end{itemize}
  \item La gasificación transforma la materia prima hidrocarbonada:
  \begin{itemize}
    \item Añade hidrógeno.
    \item Elimina carbono.
    \item Resultado: gases con alta relación H/C.
  \end{itemize}
\end{itemize}
\end{frame}
\begin{frame}{Etapas de la Gasificación de Biomasa}
\small
Un proceso típico de gasificación de biomasa incluye las siguientes etapas:

\vspace{0.5em}
\begin{enumerate}
  \item \textbf{Secado:} eliminación de humedad del material de alimentación.
  \item \textbf{Pirolisis:} descomposición térmica sin oxígeno.
  \item \textbf{Combustión parcial:} oxidación limitada de gases, vapores y char para generar calor.
  \item \textbf{Gasificación:} conversión de productos descompuestos en gases útiles.
\end{enumerate}

\vspace{0.5em}
\textbf{Resultado:} se obtiene una mezcla de gases combustibles como CO, H\textsubscript{2}, CH\textsubscript{4} y trazas de otros componentes.
\end{frame}

\begin{frame}{Gasificación vs Pirolisis y Torrefacción}
\footnotesize
\textbf{Pirolisis:}
\begin{itemize}
  \item Proceso de descomposición térmica.
  \item Se realiza en \textbf{ausencia de cualquier medio gasificante}.
  \item Genera char, líquidos y gases, dependiendo de la velocidad de calentamiento.
\end{itemize}
\textbf{Torrefacción:}
\begin{itemize}
  \item Tratamiento térmico suave de biomasa.
  \item En atmósfera inerte o con poco oxígeno.
  \item Produce un sólido mejorado (biocarbón) con mayor estabilidad y energía específica.
\end{itemize}
\textbf{Gasificación:}
\begin{itemize}
  \item Requiere un \textbf{medio gasificante}: vapor, aire u oxígeno.
  \item Reorganiza la estructura molecular para producir \textbf{gases o líquidos}.
  \item También puede \textbf{añadir hidrógeno} al producto final.
\end{itemize}
\end{frame}

\begin{frame}{Medio Gasificante: Oxígeno}
\small
\begin{itemize}
  \item El \textbf{oxígeno} es un agente gasificante común, utilizado principalmente en:
  \begin{itemize}
    \item Combustión.
    \item Gasificación parcial.
  \end{itemize}
  \item Puede suministrarse en forma pura o mediante aire.
  \item En el diagrama ternario (C–H–O), la conversión se dirige hacia la \textbf{esquina del oxígeno}.
  \begin{itemize}
    \item Baja cantidad de oxígeno: se forma principalmente CO.
    \item Alta cantidad de oxígeno: se forma CO\textsubscript{2}.
    \item Exceso de oxígeno: el proceso se convierte en \textbf{combustión}, produciendo \textbf{gas de escape} (sin poder calorífico).
  \end{itemize}
  \item Mayor formación de compuestos carbonados (CO, CO\textsubscript{2}) y menor contenido de hidrógeno.
\end{itemize}
\end{frame}

\begin{frame}{Medio Gasificante: Vapor de Agua (Steam)}

\begin{itemize}
  \item El \textbf{vapor de agua} como agente gasificante dirige la conversión hacia la \textbf{esquina del hidrógeno} en el diagrama ternario.
  \item Resultado: mayor proporción de hidrógeno en el gas producto.
  \item El gas generado tiene una \textbf{alta relación H/C}, ideal para procesos que requieren mayor contenido de H\textsubscript{2}.
  \item Su poder calorífico es \textbf{intermedio}, mayor que el generado con aire pero menor que el obtenido con oxígeno puro.
\end{itemize}
\end{frame}

\begin{frame}{Medio Gasificante: Aire}
%\small
\begin{itemize}
  \item El \textbf{aire} es el agente gasificante más económico y disponible.
  \item Contiene una gran proporción de \textbf{nitrógeno (N\textsubscript{2})}, que:
  \begin{itemize}
    \item No participa en las reacciones principales.
    \item Diluye el gas producto.
  \end{itemize}
  \item Como resultado:
  \begin{itemize}
    \item Se reduce el \textbf{poder calorífico} del gas obtenido.
    \item El gas contiene mayor volumen, pero menor energía por unidad.
    \item \textbf{Oxígeno} produce el gas con mayor poder calorífico.
    \item Le sigue el \textbf{vapor}.
    \item El \textbf{aire} genera el gas con menor poder calorífico.
      \end{itemize}
\end{itemize}
\end{frame}

\begin{frame}{Proceso de Gasificación}
\small
Un proceso típico de gasificación sigue la siguiente secuencia:

\vspace{0.5em}
\begin{enumerate}
  \item \textbf{Precalentamiento y secado} de la biomasa.
  \item \textbf{Pirolisis} y/o combustión parcial.
  \item \textbf{Gasificación del char} (reacciones de gas-sólido).
\end{enumerate}

\vspace{0.5em}
\textbf{Nota:} Aunque a menudo se modelan en serie, estas etapas no tienen límites definidos y suelen solaparse.

\vspace{0.5em}
La energía térmica para estas transformaciones proviene de reacciones exotérmicas de combustión parcial dentro del gasificador.
\end{frame}

\begin{frame}{Productos de Pirolisis y Reacciones Subsiguientes}
\small
\textbf{Durante la pirolisis:}
\begin{itemize}
  \item \textbf{Gases:} CO, H\textsubscript{2}, CH\textsubscript{4}, H\textsubscript{2}O.
  \item \textbf{Líquidos:} alquitrán, aceite, nafta.
  \item \textbf{Compuestos oxigenados:} fenoles, ácidos.
  \item \textbf{Sólido:} char.
\end{itemize}

\textbf{Reacciones posteriores:}
\begin{itemize}
  \item \textbf{Char:} participa en reacciones de gasificación, combustión y desplazamiento (shift).
  \item \textbf{Gases:} sufren craqueo, reforma, combustión parcial y reacción de desplazamiento de agua-gas.
\end{itemize}

\vspace{0.5em}
Estas interacciones generan el \textbf{gas de síntesis final:} CO, H\textsubscript{2}, CH\textsubscript{4}, CO\textsubscript{2}, H\textsubscript{2}O y otros.
\end{frame}

\begin{frame}{Poder Calorífico del Gas Producto}
\small
\textbf{Tabla – Valores caloríficos según el medio gasificante:}

\vspace{0.5em}
\begin{tabular}{|c|c|}
\hline
\textbf{Medio Gasificante} & \textbf{Poder Calorífico (MJ/Nm\textsuperscript{3})} \\
\hline
Aire & 4–7 \\
Vapor de agua & 10–18 \\
Oxígeno & 12–28 \\
\hline
\end{tabular}

\vspace{1em}
\textbf{Observaciones:}
\begin{itemize}
  \item El \textbf{oxígeno} produce el gas con mayor poder calorífico.
  \item El \textbf{vapor} genera un gas intermedio.
  \item El \textbf{aire} produce el gas con menor energía debido a la dilución con nitrógeno.
\end{itemize}
\end{frame}

\begin{frame}{Secuencia de Reacciones en la Gasificación}
\small
\textbf{Esquema de rutas potenciales:}

\vspace{0.5em}
\textbf{1. Secado:} eliminación de humedad de la biomasa.

\textbf{2. Pirolisis:} 
\begin{itemize}
  \item Gases: CO, H\textsubscript{2}, CH\textsubscript{4}, H\textsubscript{2}O.
  \item Líquidos: alquitrán, aceites, nafta.
  \item Sólido: char (carbón fijo).
  \item Compuestos oxigenados: fenoles, ácidos.
\end{itemize}

\textbf{3. Reacciones del char:}
\begin{itemize}
  \item Gasificación.
  \item Combustión parcial.
  \item Reacción de desplazamiento.
\end{itemize}

\textbf{4. Reacciones en fase gaseosa:}
\begin{itemize}
  \item Craqueo, reforma, combustión parcial, desplazamiento.
  \item Productos: CO, H\textsubscript{2}, CH\textsubscript{4}, CO\textsubscript{2}, H\textsubscript{2}O, residuos de C no convertido.
\end{itemize}
\end{frame}

\begin{frame}{Esquema del Proceso de Gasificación}
\tiny
\begin{center}
\begin{tikzpicture}[node distance=1cm, every node/.style={align=center}, >=stealth, thick]

% Etapas principales
\node[rectangle, draw, fill=blue!10] (drying) {Secado\\(Biomasa)};
\node[rectangle, draw, fill=orange!10, below=of drying] (pyrolysis) {Pirolisis\\Gases, Líquidos, Char};
\node[rectangle, draw, fill=green!10, below left=of pyrolysis, xshift=-0cm] (char) {Reacciones del Char\\Gasificación, Combustión, Shift};
\node[rectangle, draw, fill=yellow!20, below right=of pyrolysis, xshift=0cm] (gas) {Reacciones en Fase Gaseosa\\Craqueo, Reforma, Combustión, Shift};

% Conexiones
\draw[->] (drying) -- (pyrolysis);
\draw[->] (pyrolysis) -- (char);
\draw[->] (pyrolysis) -- (gas);
\draw[->, dashed] (char) -- (gas);

% Productos finales
\node[rectangle, draw, fill=cyan!15, below=of gas, yshift=-0.5cm] (products) {Productos Finales:\\CO, H\textsubscript{2}, CH\textsubscript{4}, CO\textsubscript{2}, H\textsubscript{2}O, \\Carbono no convertido};

\draw[->] (gas) -- (products);
\draw[->, dashed] (char) -- (products);

% Etiquetas de medios
\node[above left=0.2cm and 1.2cm of char] (label1) {\footnotesize\textbf{Medio gasificante:} \\O\textsubscript{2}, H\textsubscript{2}O, Aire};
\draw[->, dotted] (label1) -- (char);

\end{tikzpicture}
\end{center}
\end{frame}

\begin{frame}{Etapa de Secado en la Gasificación}
\small
\begin{itemize}
  \item La biomasa recién cortada puede contener entre \textbf{30\% y 60\%} de humedad. Algunas superan el \textbf{90\%}.
  \item Cada kg de humedad consume \textbf{≈ 2242 kJ} de energía en su evaporación — energía que no se recupera.
  \item La alta humedad reduce la eficiencia energética del gasificador.
  \item Se recomienda predesechar la humedad externa antes de alimentar la biomasa.
  \item Para un gas con buen poder calorífico, se prefiere biomasa con \textbf{10–20\%} de humedad.
  \item El secado final ocurre dentro del gasificador, donde:
  \begin{itemize}
    \item El calor de la zona caliente evapora la humedad remanente.
    \item El agua libre se elimina a temperaturas \textbf{> 100°C}.
    \item Compuestos volátiles ligeros comienzan a salir a partir de \textbf{≈ 200°C}.
  \end{itemize}
\end{itemize}
\end{frame}

\begin{frame}{Reacciones Típicas en la Gasificación (25°C)}
\scriptsize
\textbf{Reacciones del Carbono:}
\begin{itemize}
  \item R1 (Boudouard): C + CO\textsubscript{2} $\rightarrow$ 2CO \hfill $\Delta H = +172$ kJ/mol
  \item R2 (agua-gas): C + H\textsubscript{2}O $\rightarrow$ CO + H\textsubscript{2} \hfill $\Delta H = +131$ kJ/mol
  \item R3 (hidrogasificación): C + 2H\textsubscript{2} $\rightarrow$ CH\textsubscript{4} \hfill $\Delta H = –74.8$ kJ/mol
  \item R4: C + 0.5O\textsubscript{2} $\rightarrow$ CO \hfill $\Delta H = –111$ kJ/mol
\end{itemize}

\vspace{0.5em}
\textbf{Reacciones de Oxidación:}
\begin{itemize}
  \item R5: C + O\textsubscript{2} $\rightarrow$ CO\textsubscript{2} \hfill $\Delta H = –394$ kJ/mol
  \item R6: CO + 0.5O\textsubscript{2} $\rightarrow$ CO\textsubscript{2} \hfill $\Delta H = –284$ kJ/mol
  \item R7: CH\textsubscript{4} + 2O\textsubscript{2} $\rightarrow$ 2CO\textsubscript{2} + 2H\textsubscript{2}O \hfill $\Delta H = –803$ kJ/mol
  \item R8: H\textsubscript{2} + 0.5O\textsubscript{2} $\rightarrow$ H\textsubscript{2}O \hfill $\Delta H = –242$ kJ/mol
\end{itemize}
\end{frame}

\begin{frame}{Otras Reacciones en Gasificación (25°C)}
\scriptsize
\textbf{Reacción de Desplazamiento (Shift):}
\begin{itemize}
  \item R9: CO + H\textsubscript{2}O $\rightarrow$ CO\textsubscript{2} + H\textsubscript{2} \hfill $\Delta H = –41.2$ kJ/mol
\end{itemize}

\vspace{0.5em}
\textbf{Reacciones de Metanación:}
\begin{itemize}
  \item R10: 2CO + 2H\textsubscript{2} $\rightarrow$ CH\textsubscript{4} + CO\textsubscript{2} \hfill $\Delta H = –247$ kJ/mol
  \item R11: CO + 3H\textsubscript{2} $\rightarrow$ CH\textsubscript{4} + H\textsubscript{2}O \hfill $\Delta H = –206$ kJ/mol
  \item R14: CO\textsubscript{2} + 4H\textsubscript{2} $\rightarrow$ CH\textsubscript{4} + 2H\textsubscript{2}O \hfill $\Delta H = –165$ kJ/mol
\end{itemize}

\vspace{0.5em}
\textbf{Reformado con Vapor:}
\begin{itemize}
  \item R12: CH\textsubscript{4} + H\textsubscript{2}O $\rightarrow$ CO + 3H\textsubscript{2} \hfill $\Delta H = +206$ kJ/mol
  \item R13: CH\textsubscript{4} + 0.5O\textsubscript{2} $\rightarrow$ CO + 2H\textsubscript{2} \hfill $\Delta H = –36$ kJ/mol
\end{itemize}
\end{frame}

\begin{frame}{Fundamentos de la Pirolisis}
\small
\begin{itemize}
  \item La \textbf{pirolisis} no requiere un agente externo (sin aire, oxígeno o vapor).
  \item Involucra la descomposición térmica de biomasa:
  \begin{itemize}
    \item Moléculas hidrocarbonadas grandes → moléculas pequeñas.
    \item Genera gases condensables y no condensables, líquidos y char.
     \item \textbf{Pirolisis lenta / torrefacción}: se aproxima a la esquina del carbono → más char.
    \item \textbf{Pirolisis rápida}: se mueve hacia el eje C–H → más hidrocarburos líquidos.
  \end{itemize}
  \item Es una etapa previa común antes de la gasificación.
\end{itemize}
\end{frame}

\begin{frame}{Alquitrán en la Pirolisis}
\small
\begin{itemize}
  \item Uno de los productos principales de la pirolisis es el \textbf{alquitrán} (tar).
  \item Se forma por \textbf{condensación de los vapores condensables} generados.
  \item El alquitrán es un líquido pegajoso que:
  \begin{itemize}
    \item Genera serios problemas en aplicaciones industriales del gas de síntesis.
    \item Obstruye tuberías, filtros, intercambiadores y sistemas de limpieza.
  \end{itemize}
  \item existen métodos de:
  \begin{itemize}
    \item \textbf{Craqueo} del alquitrán: romper moléculas grandes en gases útiles.
    \item \textbf{Reformado}: conversión catalítica en gases no condensables como CO y H\textsubscript{2}.
  \end{itemize}
\end{itemize}
\end{frame}

\begin{frame}{Char y su Importancia en la Gasificación}
\small
\begin{itemize}
  \item La gasificación implica reacciones químicas entre hidrocarburos del combustible, vapor, CO\textsubscript{2}, O\textsubscript{2} y H\textsubscript{2}.
  \item El \textbf{char de biomasa}, generado en la pirolisis, no es carbono puro:
  \begin{itemize}
    \item Contiene compuestos hidrocarbonados con H y O.
  \end{itemize}
  \item Comparado con el coque de carbón:
  \begin{itemize}
    \item Mayor porosidad: \textbf{40–50\%} (vs. 2–18\% del carbón).
    \item Poros más grandes: \textbf{20–30 μm} (vs. ≈0.5 Å del carbón).
  \end{itemize}
  \item Es más reactivo que el char de carbón, lignito o turba.
  \item Su reactividad \textbf{aumenta con la conversión}, debido a la acción catalítica de metales alcalinos presentes.
\end{itemize}
\end{frame}

\begin{frame}{Reacciones del Char con Agentes Gasificantes}
\small
\textbf{Principales reacciones del char con agentes gasificantes:}
\vspace{0.5em}
\begin{itemize}
  \item (7.1) Char + O\textsubscript{2} $\rightarrow$ CO\textsubscript{2} y CO
  \item (7.2) Char + CO\textsubscript{2} $\rightarrow$ CO
  \item (7.3) Char + H\textsubscript{2}O $\rightarrow$ CH\textsubscript{4} y CO
  \item (7.4) Char + H\textsubscript{2} $\rightarrow$ CH\textsubscript{4}
\end{itemize}

\vspace{0.5em}
\textbf{Productos:} gases de bajo peso molecular como CO, H\textsubscript{2} y CH\textsubscript{4}

\vspace{0.5em}
\textbf{Nombres comunes:}
\begin{itemize}
  \item Reacción de Boudouard, gas-agua, hidrogasificación, etc. (ver Tabla 7.2).
\end{itemize}
\end{frame}

\begin{frame}{Naturaleza Térmica de las Reacciones del Char}
\small
\begin{itemize}
  \item Las reacciones de gasificación del char pueden ser:
  \begin{itemize}
    \item \textbf{Endotérmicas:} absorben calor.
    \item \textbf{Exotérmicas:} liberan calor.
  \end{itemize}

  \item \textbf{Reacciones endotérmicas:}
  \begin{itemize}
    \item R1: C + CO\textsubscript{2} $\rightarrow$ 2CO \hfill ($\Delta H$ = +172 kJ/mol)
    \item R2: C + H\textsubscript{2}O $\rightarrow$ CO + H\textsubscript{2} \hfill ($\Delta H$ = +131 kJ/mol)
  \end{itemize}

  \item \textbf{Reacciones exotérmicas:}
  \begin{itemize}
    \item R3: C + 2H\textsubscript{2} $\rightarrow$ CH\textsubscript{4} \hfill ($\Delta H$ = –74.8 kJ/mol)
    \item R4: C + 0.5O\textsubscript{2} $\rightarrow$ CO \hfill ($\Delta H$ = –111 kJ/mol)
    \item R5: C + O\textsubscript{2} $\rightarrow$ CO\textsubscript{2} \hfill ($\Delta H$ = –394 kJ/mol)
  \end{itemize}

  \item Los valores de entalpía están referidos a \textbf{25°C}.
\end{itemize}
\end{frame}

\begin{frame}{Velocidad de las Reacciones del Char}
\small
\begin{itemize}
  \item Las reacciones químicas del char (C) tienen velocidades finitas y dependen de:
  \begin{itemize}
    \item La \textbf{reactividad del char}.
    \item El \textbf{potencial del medio gasificante}.
  \end{itemize}

  \item \textbf{Velocidades relativas} de reacciones a 800°C (Walker et al., 1959):
  \begin{itemize}
    \item C + 0.5O\textsubscript{2} $\rightarrow$ CO \hfill $\sim 10^5$ (más rápida)
    \item C + H\textsubscript{2}O $\rightarrow$ CO + H\textsubscript{2} \hfill $\sim 10^3$
    \item C + CO\textsubscript{2} $\rightarrow$ 2CO \hfill $\sim 10^1$
    \item C + 2H\textsubscript{2} $\rightarrow$ CH\textsubscript{4} \hfill $\sim 3 \times 10^{-3}$ (más lenta)
  \end{itemize}

  \item Relación general:
  \[
  R_{C+O_2} \gg R_{C+H_2O} \gg R_{C+CO_2} \gg R_{C+H_2}
  \]
\end{itemize}
\end{frame}

\begin{frame}{Comparación de Reacciones del Char con Diferentes Agentes}
\small
\textbf{Orden de reactividad de agentes gasificantes:}
\[
\text{Oxígeno} > \text{Vapor de agua} > \text{Dióxido de carbono} > \text{Hidrógeno}
\]

\vspace{0.5em}
\textbf{Productos formados según el medio:}
\begin{itemize}
  \item \textbf{Oxígeno (O\textsubscript{2})} $\rightarrow$ CO (rápida, exotérmica)
  \item \textbf{Vapor (H\textsubscript{2}O)} $\rightarrow$ CO + H\textsubscript{2} (reacción gas-agua)
  \item \textbf{CO\textsubscript{2}} $\rightarrow$ 2CO (reacción de Boudouard)
  \item \textbf{Hidrógeno (H\textsubscript{2})} $\rightarrow$ CH\textsubscript{4} (hidrogasificación, muy lenta)
\end{itemize}

\vspace{0.5em}
\textbf{Notas:}
\begin{itemize}
  \item La reacción con vapor puede también generar CH\textsubscript{4} y CO\textsubscript{2} bajo ciertas condiciones.
  \item La reacción con O\textsubscript{2} es tan rápida que consume el oxígeno antes de que participen otras reacciones.
\end{itemize}
\end{frame}

\begin{frame}{Reacción de Boudouard}
\small
\begin{itemize}
  \item Es la gasificación del char (C) con dióxido de carbono:
  \[
  \text{C} + \text{CO}_2 \rightarrow 2\text{CO} \quad \text{(Reacción R1)}
  \]
  \item Reacción endotérmica (+172 kJ/mol).
  \item Relevante en procesos de gasificación donde el CO\textsubscript{2} actúa como agente gasificante.
  \item Velocidad de reacción \textbf{muy baja por debajo de 1000 K}.
  \item Importancia:
  \begin{itemize}
    \item Contribuye a la generación de CO.
    \item Influyente en la composición del gas de síntesis.
  \end{itemize}
\end{itemize}
\end{frame}

\begin{frame}{Mecanismo de la Reacción de Boudouard}
\small
\textbf{Según Blasi (2009), la reacción se desarrolla en tres pasos:}

\vspace{0.5em}
\begin{enumerate}
  \item \textbf{Paso 1:} Activación del sitio libre de carbono (C\textsubscript{fas}):
  \[
  \text{C}_{\text{fas}} + \text{CO}_2 \xrightarrow{k_{b1}} \text{C(O)} + \text{CO}
  \]
  \item \textbf{Paso 2:} Reversibilidad – recombinación con CO:
  \[
  \text{C(O)} + \text{CO} \xrightarrow{k_{b2}} \text{C}_{\text{fas}} + \text{CO}_2
  \]
  \item \textbf{Paso 3:} Descomposición del complejo C–O:
  \[
  \text{C(O)} \xrightarrow{k_{b3}} \text{CO}
  \]
\end{enumerate}

\vspace{0.5em}
\textbf{Notas:}
\begin{itemize}
  \item El complejo superficial C(O) es intermedio clave.
  \item La reacción puede avanzar o revertirse según las condiciones.
\end{itemize}
\end{frame}

\begin{frame}{Reacción Gas–Agua (Water–Gas)}
\small
\begin{itemize}
  \item Es una de las reacciones más importantes de la gasificación del char.
  \[
  \text{C} + \text{H}_2\text{O} \rightarrow \text{CO} + \text{H}_2 \quad \text{(R2)}
  \]
  \item Reacción endotérmica: $\Delta H = +131$ kJ/mol.
  \item Produce monóxido de carbono y gas hidrógeno.
  \item Es más rápida que la reacción de Boudouard, pero más lenta que la reacción con oxígeno.
  \item Clave para procesos que buscan maximizar la producción de H\textsubscript{2}.
\end{itemize}
\end{frame}

\begin{frame}{Mecanismo de la Reacción Gas–Agua (Blasi, 2009)}
\small
\textbf{Etapas del mecanismo sobre un sitio activo libre de carbono (C\textsubscript{fas}):}

\vspace{0.5em}
\begin{enumerate}
  \item \textbf{Paso 1:} Disociación del H\textsubscript{2}O:
  \[
  \text{C}_{\text{fas}} + \text{H}_2\text{O} \xrightarrow{k_{w1}} \text{C(O)} + \text{H}_2
  \]
  \item \textbf{Paso 2 (reversible):} Recombinación de C(O) con H\textsubscript{2}:
  \[
  \text{C(O)} + \text{H}_2 \xrightarrow{k_{w2}} \text{C}_{\text{fas}} + \text{H}_2\text{O}
  \]
  \item \textbf{Paso 3:} Formación de monóxido de carbono:
  \[
  \text{C(O)} \xrightarrow{k_{w3}} \text{CO}
  \]
\end{enumerate}

\vspace{0.5em}
\textbf{Producto global:} CO + H\textsubscript{2}
\end{frame}

\begin{frame}{Efecto Inhibidor del Hidrógeno}
\small
\begin{itemize}
  \item Algunos modelos (Blasi, 2009) consideran la formación de complejos inhibidores:
  \begin{align*}
  \text{C}_{\text{fas}} + \text{H}_2 &\rightarrow \text{C(H)}_2 \quad \text{(7.14)} \\
  \text{C}_{\text{fas}} + 0.5\text{H}_2 &\rightarrow \text{C(H)} \quad \text{(7.15)}
  \end{align*}

  \item Estos complejos \textbf{bloquean los sitios activos}, reduciendo la velocidad de gasificación.

  \item Ejemplo (Barrio et al., 2001):
  \begin{itemize}
    \item Una atmósfera con \textbf{30\% de H\textsubscript{2}} puede reducir la velocidad de gasificación hasta \textbf{15 veces}.
  \end{itemize}

  \item \textbf{Estrategia para acelerar la reacción:}
  \begin{itemize}
    \item Eliminar continuamente el H\textsubscript{2} del sitio de reacción.
  \end{itemize}
\end{itemize}
\end{frame}

\begin{frame}{Reacción Shift (Water–Gas Shift)}
\small
\begin{itemize}
  \item Es una reacción en fase gaseosa entre productos intermedios:
  \[
  \text{CO} + \text{H}_2\text{O} \rightarrow \text{CO}_2 + \text{H}_2 \quad (\Delta H = -41.2\ \text{kJ/mol})
  \]
  \item Aumenta el contenido de \textbf{hidrógeno} a costa del monóxido de carbono.
  \item También llamada “reacción water–gas shift”, distinta de la reacción gas–agua (R2).
  \item Clave en la etapa posterior del gasificador para producir gas de síntesis (syngas) con proporción H\textsubscript{2}/CO adecuada.
\end{itemize}
\end{frame}

\begin{frame}{Equilibrio y Condiciones Operativas de la Reacción Shift}
\small
\begin{itemize}
  \item Reacción \textbf{ligeramente exotérmica}:
  \[
  \text{Mayor temperatura} \Rightarrow \text{menor rendimiento de H\textsubscript{2}}, pero \text{mayor velocidad}.
  \]
  \item \textbf{No sensible a la presión.}
  \item A \textbf{T > 1000°C}, alcanza el equilibrio rápidamente.
  \item A \textbf{T baja (180–270°C)}, necesita catalizadores heterogéneos:
  \begin{itemize}
    \item Óxidos de cobre–zinc–aluminio.
    \item Óxidos de hierro promocionados con cromo.
  \end{itemize}
  \item Catalizadores comerciales:
  \begin{itemize}
    \item Cu a 300–510°C.
    \item Cu/Zn/Al\textsubscript{2}O\textsubscript{3} a 180–270°C.
  \end{itemize}
  \item \textbf{Óptimo de rendimiento:} alrededor de \textbf{225°C}.
\end{itemize}
\end{frame}

\begin{frame}{Reacción de Hidrogasificación (R3)}
\small
\begin{itemize}
  \item Se lleva a cabo en presencia de hidrógeno:
  \[
  \text{C} + 2\text{H}_2 \rightarrow \text{CH}_4
  \]
  \item Reacción muy lenta comparada con otras.
  \item Relevante solo cuando se busca producir \textbf{gas natural sintético}.
  \item No tiene gran impacto en la gasificación general convencional.
\end{itemize}
\end{frame}

\begin{frame}{Reacciones de Combustión del Char}
\small
\begin{itemize}
  \item La mayoría de las reacciones de gasificación son \textbf{endotérmicas}.
  \item Se permite una cantidad controlada de \textbf{combustión exotérmica} para:
  \begin{itemize}
    \item Suministrar calor para: secado, pirolisis y reacciones endotérmicas.
  \end{itemize}
  \item Principales reacciones de combustión del char:
  \begin{align*}
  \text{R5: } & \text{C + O}_2 \rightarrow \text{CO}_2 \quad \Delta H = -394\ \text{kJ/mol} \\
  \text{R4: } & \text{C + 0.5 O}_2 \rightarrow \text{CO} \quad \Delta H = -111\ \text{kJ/mol}
  \end{align*}
  \item R5 entrega más calor, pero R4 produce CO como gas útil.
\end{itemize}
\end{frame}

\begin{frame}{Coeficiente de Partición y Velocidades}
\small
\textbf{Reacciones R4 y R5 pueden ocurrir simultáneamente.}

\begin{itemize}
  \item El \textbf{coeficiente de partición} $\beta$ indica cómo se reparte el oxígeno:
  \[
  \beta \in [1, 2], \quad \text{depende de la temperatura (T)}
  \]
  \item Forma combinada de las reacciones:
  \[
  \beta\,\text{C} + \text{O}_2 \rightarrow 2(\beta-1)\,\text{CO} + (2-\beta)\,\text{CO}_2
  \]
  \item Ecuación de Arthur (1951) para $\beta$:
  \[
  \beta = \frac{[\text{CO}]}{[\text{CO}_2]} = 52400\,e^{-6234/T}
  \]
  \item Las reacciones de \textbf{combustión son mucho más rápidas} que las de gasificación.
\end{itemize}
\end{frame}

\begin{frame}{Efecto del Tamaño de Partícula y Comparación de Velocidades}
\small
\textbf{Comparación entre combustión y gasificación del char a 900°C:}

\vspace{0.5em}
\begin{tabular}{|c|c|c|c|}
\hline
\textbf{Tamaño (μm)} & \textbf{Comb. (min⁻¹)} & \textbf{Gasif. (min⁻¹)} & \textbf{Relación} \\
\hline
6350 & 0.648 & 0.042 & 15.4 \\
841  & 5.04  & 0.317 & 15.9 \\
74   & 55.9  & 0.975 & 57.3 \\
\hline
\end{tabular}

\vspace{0.5em}
\textbf{Conclusiones:}
\begin{itemize}
  \item Partículas más pequeñas: mayor velocidad de reacción.
  \item En combustión: la temperatura del char puede superar la del lecho.
  \item En gasificación: la temperatura del char ≈ temperatura del lecho (por equilibrio térmico).
\end{itemize}
\end{frame}


\end{document}